% GENERATED by scripts/build-lecture-slides.py; do not edit. Portable profile: XeTeX, theme and fonts from the core.
\makeatletter\IfFileExists{core/beamer/beamerthememx.sty}{\def\mxcorepath{core/}\def\input@path{{core/beamer/}{core/}}}{\def\mxcorepath{../core/}\def\input@path{{../core/beamer/}{../core/}}}\makeatother
\documentclass[aspectratio=169,10pt,t]{beamer}
\usetheme{mx}
\usepackage{graphicx,booktabs}
\usepackage[french]{babel}
\setlength{\parskip}{.45em}
\hypersetup{pdfauthor={},pdfsubject={Lecture slides},pdfkeywords={},pdftitle={Fonctions exponentielle et logarithme}}
\providecommand{\N}{\mathbb{N}}
\providecommand{\Z}{\mathbb{Z}}
\providecommand{\D}{\mathbb{D}}
\providecommand{\Q}{\mathbb{Q}}
\providecommand{\R}{\mathbb{R}}
\providecommand{\C}{\mathbb{C}}
\providecommand{\vect}[1]{\overrightarrow{#1}}
\providecommand{\vi}{\vec{\imath}}
\providecommand{\vj}{\vec{\jmath}}
\providecommand{\vk}{\vec{k}}
\providecommand{\norm}[1]{\left\lVert #1\right\rVert}
\providecommand{\vc}[2]{\begin{pmatrix}#1\\#2\end{pmatrix}}
\providecommand{\deter}[2]{\det\left(#1,#2\right)}
\providecommand{\abs}[1]{\left\lvert #1\right\rvert}
\providecommand{\dd}{\,\mathrm{d}}
\providecommand{\e}{\mathrm{e}}
\providecommand{\Card}{\operatorname{Card}}
\providecommand{\E}{\operatorname{E}}
\providecommand{\V}{\operatorname{V}}
\providecommand{\degre}{^{\circ}}
\providecommand{\un}[1]{\,\mathrm{#1}}
\providecommand{\sci}[2]{#1\times10^{#2}}
\providecommand{\pgcd}{\operatorname{PGCD}}
\providecommand{\ppcm}{\operatorname{PPCM}}
\title{Fonctions exponentielle et logarithme}
\subtitle{Cours et exemples}
\date{}
\author{}
\begin{document}
\mxtitleframe
\begin{frame}{Fonction exponentielle}
La \emph{fonction exponentielle} est l'unique fonction \(f\) définie et dérivable sur \(\mathbb{R}\) et vérifiant : 
                  \[   f' = f \qquad \text{ et } \qquad f(0) = 1   \,. \]
            On note la fonction par \(\mathrm{exp}\).
\end{frame}
\begin{frame}{Fonction exponentielle}
La fonction exponentielle peut en fait être considérée comme une fonction puissance.
            Il existe un nombre réel \(e\) (\(\simeq 2,718\,281\,828\,459\,045\)) 
            tel que \(\exp{x} = e^{x}\) pour tout \(x\in\mathbb{Q}\).
            
            Ainsi, dorénavant, on considèrera la fonction exponentielle comme 
            une généralisation d'une fonction puissance.
            On dénotera alors \(\exp{x}\) par \(\mathrm{e}^{x}\).
\end{frame}
\begin{frame}{Fonction exponentielle}
Soient \(a\) et \(b\) deux nombres réels.
\end{frame}
\begin{frame}{Fonction exponentielle (2)}
\begin{itemize}

                  \item \(\mathrm{e}^{0} = 1\)
                  
\item \(\mathrm{e}^{a+b} = \mathrm{e}^{a} \cdot \mathrm{e}^{b}\)
                  
\item \(\mathrm{e}^{-a} = \dfrac{1}{\mathrm{e}^{a}}\)
                  
\item \(\mathrm{e}^{a-b} = \dfrac{\mathrm{e}^{a}}{\mathrm{e}^{b}}\)
                  
\item \(\forall n \in \mathbb{Q},\ \left(\mathrm{e}^{a}\right)^n = \mathrm{e}^{na}\).

\end{itemize}
\end{frame}
\begin{frame}{Fonction exponentielle}
La fonction \(\exp\) est définie sur \(\mathbb{R}\).
            De plus, \(\mathrm{e}^{x} > 0\) pour tout \(x\in\mathbb{R}\) (car \(\mathrm{e}^{x} = \left(\mathrm{e}^{x/2}\right)^2\) et \(\mathrm{e}^{x/2}\,\mathrm{e}^{-x/2} = 1\), donc \(\mathrm{e}^{x/2} \neq 0\)).
            Comme \(\exp' x = \exp x > 0\), alors la fonction \(\exp\) est strictement croissante sur \(\mathbb{R}\).
            Il suit que, pour tout \((x,y)\in\mathbb{R}^2\) :
                  \[
                        x<y \Rightarrow \mathrm{e}^{x} < \mathrm{e}^{y}      \;.
                  \]
\end{frame}
\begin{frame}{Fonction exponentielle}
\[\begin{aligned}
                  \lim_{x\rightarrow-\infty} \mathrm{e}^{x} & = 0     & \lim_{x\rightarrow+\infty} \mathrm{e}^{x} & = +\infty \\
                  \lim_{x\rightarrow-\infty} x \mathrm{e}^{x} & = 0   & \lim_{x\rightarrow+\infty} \dfrac{\mathrm{e}^{x}}{x} & = +\infty \\
                \\ \text{ou plus généralement, pour tout $n\in\mathbb{N}$ (<< croissances comparées >> : en $+\infty$, l'exponentielle l'emporte sur toute puissance de $x$)} \\
                  \lim_{x\rightarrow-\infty} x^n \mathrm{e}^{x} & = 0 & \lim_{x\rightarrow+\infty} \dfrac{\mathrm{e}^{x}}{x^n} & = +\infty\;.
            \end{aligned}\]
            On a aussi 
                  \[
                        \lim_{x\rightarrow 0} \dfrac{\mathrm{e}^{x}-1}{x} = 1      \;.
                  \]
\end{frame}
\begin{frame}{Fonction exponentielle}
\textbf{Fonctions composées avec la fonction exponentielle}

 Il suit de la formule de dérivation d'une fonction composée que 
            \(\left( \mathrm{e}^{u} \right)' = u' \mathrm{e}^{u}\).
            Ainsi, \(u\) et \(\mathrm{e}^{u}\) ont même sens de variations.
            Par exemple, pour un réel \(k\), \(\left(\mathrm{e}^{kx}\right)' = k\,\mathrm{e}^{kx}\).
\end{frame}
\begin{frame}{Fonction exponentielle}
Étude de la fonction \(\exp\)

Posons \(f(x) = \mathrm{e}^{x}\).
\end{frame}
\begin{frame}{Fonction exponentielle (2)}
\begin{itemize}
\item \textbf{Domaine de définition :}  La fonction \(\exp\) est définie sur \(\mathbb{R} = ]-\infty,+\infty[\).
            
\item \textbf{Limites aux bornes :}  \(\lim_{x\rightarrow -\infty} \mathrm{e}^{x} = 0\) et 
                                        \(\lim_{x\rightarrow +\infty} \mathrm{e}^{x} = +\infty\)
            
\item \textbf{Dérivée :}  Pour tout \(x\in\mathbb{R}\), \(f'(x) = \mathrm{e}^{x}\).

\end{itemize}
\end{frame}
\begin{frame}{Fonction exponentielle (3)}
\begin{itemize}
\item \textbf{Tableau des variations :}  

                  table[H]

                  Tableau des variations de la fonction exponentielle.
                  
                  table

\end{itemize}
\end{frame}
\begin{frame}{Fonction exponentielle (4)}
\begin{itemize}
\item \textbf{Valeurs remarquables :}  \(f(0) = 1\) donc la courbe passe par le point de coordonnées \((0,1)\).
                                          De plus, la dérivée en ce point est égale à \(1\) :
                                          la tangente en ce point a pour équation \(y = x + 1\),
                                          elle est parallèle à la première bissectrice.
                                          \(f(1) = \mathrm{e}^{}\) donc la courbe passe par le point de coordonnées \((1,\mathrm{e}^{})\). 

\end{itemize}
\end{frame}
\begin{frame}{Fonction exponentielle (5)}
\begin{itemize}
\item \textbf{Courbe représentative :} 
                  Avant de tracer la courbe représentative de \(f\), étudions d'abord ses branches infinies.
                  
\begin{itemize}

                        \item Comme \(\lim_{x\rightarrow-\infty} f(x) = 0\), 
                              alors la courbe admet la droite d'équation \(y=0\) comme asymptote horizontale
                              au voisinage de \(-\infty\).
                        
\item Comme \(\lim_{x\rightarrow+\infty} f(x) = +\infty\) et que
                              \(\lim_{x\rightarrow+\infty} \dfrac{f(x)}{x} = +\infty\) alors
                              la courbe admet une branche parabolique de direction \((y'Oy)\) au voisinage de \(+\infty\).

\end{itemize}

\end{itemize}
\end{frame}
\begin{frame}{Fonction Logarithme}
On a vu que la fonction exponentielle est continue (elle est même dérivable)
            et strictement croissante sur \(\mathbb{R}\).
            Elle définit donc une bijection de \(\mathbb{R}\) vers \(\mathbb{R}^*_+\).
            On définit la fonction logarithme népérien comme étant la fonction réciproque de la fonction exponentielle.
            Ainsi, 
            \(\ln : \mathbb{R}^*_+   \rightarrow   \mathbb{R}\) vérifie :
                  \[
                        \left\lbrace
                              \begin{array}{ll}
                                    \forall x\in\mathbb{R}      ,&  \ln(\mathrm{e}^{x}) = x \\
                                    \forall x\in\mathbb{R}^*_+  ,&  \mathrm{e}^{\ln x} = x  \;.
                              \end{array}
                        \right.
                  \]
\end{frame}
\begin{frame}{Fonction Logarithme}
Ainsi, dans un repère orthonormé, la courbe de la fonction \(\exp\) et celle de la fonction \(\ln\) sont 
      symétriques par rapport à la première bissectrice (droite d'équation \(y = x\)).
\end{frame}
\begin{frame}{Fonction Logarithme}
\begin{itemize}

                  \item \(\ln 1 = 0\) car \(\mathrm{e}^{0} = 1\)
                  
\item \(\ln \mathrm{e}^{} = 1\) car \(\mathrm{e}^{1} = \mathrm{e}^{}\).

\end{itemize}
\end{frame}
\begin{frame}{Fonction Logarithme (2)}
Comme la fonction exponentielle est une bijection continue et strictement croissante de \(\mathbb{R}\) sur \(]0,+\infty[\), 
            la fonction logarithme népérien est une bijection continue et strictement croissante de \(]0,+\infty[\) sur \(\mathbb{R}\).
            En particulier, \(\ln x < 0\) pour \(0 < x < 1\) et \(\ln x > 0\) pour \(x > 1\).
\end{frame}
\begin{frame}{Fonction Logarithme}
Soient \(a\), \(b\) et \(x\) des nombres réels strictement positifs, \(\lambda\) un réel et \(n\) un entier relatif.
            On a alors :

            \[\begin{aligned}
                        \ln (a\cdot b) & = \ln a + \ln b                                                &       
                        \ln (a^n) & = n \cdot \ln a                                                     \\
                        \ln \dfrac{1}{a} & = - \ln a                                                     & 
                        \ln \dfrac{a}{b} & = \ln a - \ln b                                               \\
                        \ln a < \ln b & \iff a < b                                                      & 
                        \ln a = \ln b & \iff a = b                                                      \\
                        \ln x = \lambda & \iff x = \mathrm{e}^{\lambda}      & 
                        \ln x < \lambda & \iff x < \mathrm{e}^{\lambda}      
            \end{aligned}\]
            et, pour tout réel \(t\) et tout réel \(\mu > 0\) :
            \[\begin{aligned}
                        \mathrm{e}^{t} = \mu & \iff t = \ln{\mu} &
                        \mathrm{e}^{t} < \mu & \iff t < \ln{\mu} \;.
            \end{aligned}\]
\end{frame}
\begin{frame}{Fonction Logarithme}
La fonction logarithme népérien est dérivable sur \(]0,+\infty[\) avec
                  \[
                        \left( \ln x \right)' = \dfrac{1}{x}      \;.
                  \]
\end{frame}
\begin{frame}{Fonction Logarithme}
Posons \(f(x) = \mathrm{e}^{x}\). On sait alors que \(f^{-1} (x) = \ln x\).
                  En utilisant la formule de dérivation de la fonction réciproque, on a :
                        \[\begin{aligned}
                              \left(f^{-1}\right)' (x)  & = \dfrac{1}{f' \left( f^{-1} (x) \right) }      \\
                              \left(\ln (x) \right)'    & = \dfrac{1}{\mathrm{e}^{\ln x}}                       \\
                              \left(\ln (x) \right)'    & = \dfrac{1}{x}                                  \,.
                        \end{aligned}\]
\end{frame}
\begin{frame}{Fonction Logarithme}
\[\begin{aligned}
                        \lim_{x\rightarrow 0^+} \ln x                & = -\infty            &
                        \lim_{x\rightarrow 0^+} x\cdot\ln x          & = 0                  \\
                        \lim_{x\rightarrow +\infty} \ln x            & = +\infty            &
                        \lim_{x\rightarrow +\infty} \dfrac{\ln x}{x}  & = 0                  
                  \end{aligned}\]
                  Plus généralement (croissances comparées), pour tout entier \(n \geq 1\) :
                  \(\displaystyle \lim_{x\rightarrow +\infty} \dfrac{\ln x}{x^n} = 0\) et \(\displaystyle \lim_{x\rightarrow 0^+} x^n \ln x = 0\).
                  \[
                        \lim_{x\rightarrow 0} \dfrac{\ln (1 + x)}{x} = 1      \;.
                  \]
\end{frame}
\begin{frame}{Fonction Logarithme}
\textbf{Fonctions composées avec la fonction logarithme népérien}

 Soit \(u\) une fonction dérivable sur un intervalle \(I\) tel que \(u(x)>0\) pour tout \(x\in I\).
            Il suit de la formule de dérivation d'une fonction composée que 
            \(\left( \ln {u} \right)' = \dfrac{u'}{u}\).
            Ainsi, \(u\) et \(\ln{u}\) ont même sens de variations.
            Si \(u\) est dérivable et ne s'annule pas sur \(I\), on a de même
            \(\left( \ln |u| \right)' = \dfrac{u'}{u}\) ; en particulier \(\left(\ln|x|\right)' = \dfrac{1}{x}\) sur \(]-\infty,0[\) et sur \(]0,+\infty[\).
\end{frame}
\begin{frame}{Fonction Logarithme}
On définit la fonction logarithme décimal par :
                        \[
                              \begin{array}{rcl}
                                    \log :\ ]0,+\infty[ & \rightarrow & \mathbb{R}      \\
                                                     x & \mapsto     & \dfrac{\ln x}{\ln 10}
                              \end{array}
                        \]
\end{frame}
\begin{frame}{Fonction Logarithme}
Si la fonction \(\ln\) est la réciproque de la fonction \(x\mapsto \mathrm{e}^{x}\),
                  la fonction \(\log\) peut être considérée comme étant la réciproque de la fonction \(x\mapsto 10^x\).
                  
                  On a \(\log 10 = 1\), \(\log(10^p) = p\) pour tout entier \(p\), et \(\log\) vérifie les mêmes propriétés algébriques que \(\ln\) ;
                  sa dérivée est \(x \mapsto \dfrac{1}{x \ln 10}\).
                  Si \(n \geq 1\) est un nombre entier, l'écriture décimale de \(n\) comporte 
                  \(\mathrm{E}\!\left(\log (n)\right)+1\) chiffres (\(\mathrm{E}\) désigne la partie entière) :
                  en effet, si \(n\) a \(p\) chiffres, \(10^{p-1} \leq n < 10^p\), donc \(p - 1 \leq \log n < p\).
                        example
                              Le nombre \(2^{77\, 232\, 917} - 1\) est un nombre premier, découvert en décembre 2017 par le projet de calcul GIMPS.
                              Comme \(2^{77\, 232\, 917}\) n'est pas une puissance de \(10\) (il n'est pas divisible par \(5\)),
                              \(2^{77\, 232\, 917} - 1\) comporte le même nombre de chiffres que \(2^{77\, 232\, 917}\).
                              Pour calculer le nombre de chiffres qu'il comporte, on effectue :
                                    \[
                                           \log (2^{77\, 232\, 917})
                                    =      \dfrac{\ln (2^{77\, 232\, 917})}{\ln 10} 
                                    =      \dfrac{77\, 232\, 917 \cdot \ln 2}{\ln 10}
                                    \simeq 23\,249\,424{,}67
                                    .
                                    \]
                              Ainsi, ce nombre comporte \(23\,249\,425\) chiffres.
                        example
\end{frame}
\begin{frame}{Fonction Logarithme}
Soit \(a\) un nombre réel strictement positif et différent de \(1\) (pour que \(\ln a \neq 0\)).
                  On définit la fonction logarithme de base \(a\) par :
                        \[
                              \begin{array}{rcl}
                                    \log_a :\ ]0,+\infty[ & \rightarrow & \mathbb{R}                        \\
                                                        x & \mapsto     & \dfrac{\ln x}{\ln a}               \;.
                              \end{array}
                        \]
\end{frame}
\begin{frame}{Fonction Logarithme}
On a \(\log_a 1 = 0\), \(\log_a a = 1\), et \(\log_{10} = \log\).
            La fonction \(\log_a\) est croissante si \(a > 1\), décroissante si \(0 < a < 1\) ; c'est la réciproque de la fonction \(x \mapsto a^x\) définie ci-dessous.
\end{frame}
\begin{frame}{Fonction Logarithme}
Soit \(a\) un réel strictement positif. Pour tout réel \(x\), on pose
                  \[ a^x = \mathrm{e}^{x \ln a} \;. \]
                  Lorsque \(x\) est un entier, on retrouve la puissance habituelle (car \(\mathrm{e}^{n \ln a} = \left(\mathrm{e}^{\ln a}\right)^n = a^n\)).
\end{frame}
\begin{frame}{Fonction Logarithme}
Soient \(a > 0\), \(b > 0\) et \(x\), \(y\) des réels. Alors
                  \(a^{x+y} = a^x a^y\), \(a^{-x} = \dfrac{1}{a^x}\), \(\left(a^x\right)^y = a^{xy}\), \((ab)^x = a^x b^x\) et \(\ln\left(a^x\right) = x \ln a\).

La fonction \(x \mapsto a^x\) est dérivable sur \(\mathbb{R}\) et \(\left(a^x\right)' = (\ln a)\, a^x\).
\end{frame}
\begin{frame}{Fonction Logarithme (2)}
\begin{itemize}

                        \item Si \(a > 1\), elle est strictement croissante, \(\lim_{x \rightarrow -\infty} a^x = 0\) et \(\lim_{x \rightarrow +\infty} a^x = +\infty\) ;
                        
\item si \(0 < a < 1\), elle est strictement décroissante, \(\lim_{x \rightarrow -\infty} a^x = +\infty\) et \(\lim_{x \rightarrow +\infty} a^x = 0\) ;
                        
\item si \(a = 1\), elle est constante, égale à \(1\).

\end{itemize}
\end{frame}
\begin{frame}{Fonction Logarithme (3)}
Pour \(a \neq 1\), \(y = a^x \iff x = \log_a y\) (avec \(y > 0\)).
\end{frame}
\begin{frame}{Fonction Logarithme}
Soit \(\alpha\) un réel. La fonction \(x \mapsto x^\alpha = \mathrm{e}^{\alpha \ln x}\) est définie et dérivable sur \(]0,+\infty[\) et
                  \(\left(x^\alpha\right)' = \alpha\, x^{\alpha - 1}\).
                  Par exemple, \(\sqrt{x} = x^{1/2}\) pour \(x > 0\).
\end{frame}
\end{document}
